\notepage{50}{Normalized spherical harmonics}
Define the highest-weight angular function by
$Y_\ell^\ell(\theta,\phi)\propto\sin^\ell\theta\ee^{\ii\ell\phi}$.
With
\[
 \dd^3r=r^2\dd r\,\dd\Omega,\qquad
 \dd\Omega=\sin\theta\,\dd\theta\,\dd\phi,
\]
choose $\int\abs{Y_\ell^\ell}^2\dd\Omega=1$.
The needed polar integral is
\[
 \int_0^\pi\sin^{2\ell+1}\theta\,\dd\theta
 =\frac{2^{2\ell+1}(\ell!)^2}{(2\ell+1)!}.
\]
Using the conventional Condon--Shortley phase gives
\[
 \boxed{Y_\ell^\ell(\theta,\phi)
 =\frac{(-1)^\ell}{2^\ell\ell!}
   \sqrt{\frac{(2\ell+1)!}{4\pi}}
   \sin^\ell\theta\,\ee^{\ii\ell\phi}.}
\]

\smallheading{Generating the other magnetic states}
For $\ell>0$, the first lowering step gives
\[
 L_-\ket{\ell,\ell}=\hbar\sqrt{2\ell}\ket{\ell,\ell-1}.
\]
Hence, using the differential operator from note 49,
\[
 Y_\ell^{\ell-1}
 =-\frac{\ee^{-\ii\phi}}{\sqrt{2\ell}}
   \left(\partial_\theta-\ii\cot\theta\partial_\phi\right)
   Y_\ell^\ell.
\]
Before including its common normalization and phase, the differential
operation on the highest-weight angular factor is
\begin{align*}
 &\left(\partial_\theta-\ii\cot\theta\partial_\phi\right)
   \left(\sin^\ell\theta\ee^{\ii\ell\phi}\right)\\
 &\qquad=\left(\ell\cos\theta\sin^{\ell-1}\theta
       +\ell\cos\theta\sin^{\ell-1}\theta\right)\ee^{\ii\ell\phi}\\
 &\qquad=2\ell\cos\theta\sin^{\ell-1}\theta\ee^{\ii\ell\phi}.
\end{align*}
Repeated lowering generates all $Y_\ell^m$, called the spherical
harmonics. One may also start at $m=-\ell$ and raise toward $m=0$.
\[
 \ket{\ell,\ell}\xrightarrow{\ L_-\ }\cdots
 \xrightarrow{\ L_-\ }\ket{\ell,0}
 \xleftarrow{\ L_+\ }\cdots\xleftarrow{\ L_+\ }\ket{\ell,-\ell}.
\]
The corresponding phase relation is
\[
 \boxed{Y_\ell^{-m}(\theta,\phi)=(-1)^m
   \bigl(Y_\ell^m(\theta,\phi)\bigr)^*.}
\]
\editorial{The minus sign in the differential form of $L_-$ is kept
explicitly; it is required for consistency with the chosen
Condon--Shortley phase and the positive ladder coefficients.}
